% figA5_crossdomain.tex — Figure A5: Cross-domain mechanistic evidence
% % figA5_crossdomain.tex — Figure A5: Cross-domain mechanistic evidence
% % figA5_crossdomain.tex — Figure A5: Cross-domain mechanistic evidence
% % figA5_crossdomain.tex — Figure A5: Cross-domain mechanistic evidence
% \input{figures/figA5_crossdomain}

\begin{figure*}[t]
\centering
\begin{subfigure}[b]{0.31\textwidth}
\begin{tikzpicture}
\begin{axis}[
  pmh bar,
  symbolic x coords={T01,T02},
  xtick=data,
  xticklabels={T01 (CIFAR-10), T02 (Graph)},
  xticklabel style={font=\footnotesize, rotate=12, anchor=east},
  ylabel={Embedding drift at $\sigma{=}0.1$},
  ymin=0, ymax=0.82,
  legend pos=north east,
  width=0.88\linewidth, height=5.5cm,
  bar width=4.5pt,
]
\addplot[fill=colB0,   draw=none] coordinates {(T01,0.692)(T02,0.375)};
\addlegendentry{B0}
\addplot[fill=colVAT,  draw=none] coordinates {(T01,0.699)(T02,0.304)};
\addlegendentry{VAT}
\addplot[fill=colPMH,  draw=none] coordinates {(T01,0.385)(T02,0.021)};
\addlegendentry{E1 (PMH)}
\node[font=\tiny, color=colPMH] at (axis cs:T02,0.12) {18$\times$ lower};
\end{axis}
\end{tikzpicture}
\caption{Embedding drift.}
\end{subfigure}
\hfill
\begin{subfigure}[b]{0.34\textwidth}
\begin{tikzpicture}
\begin{axis}[
  pmh bar,
  symbolic x coords={S1,S2,S3,S4},
  xtick=data,
  xlabel={ResNet stage},
  ylabel={Stage drift (lower $=$ better)},
  ymin=0, ymax=14.5,
  legend pos=north west,
  width=0.88\linewidth, height=5.5cm,
  bar width=3.5pt,
]
\addplot[fill=colB0,    draw=none]
  coordinates {(S1,1.82)(S2,1.59)(S3,1.92)(S4,12.68)};
\addlegendentry{B0}
\addplot[fill=colVAT,   draw=none]
  coordinates {(S1,1.67)(S2,0.82)(S3,0.75)(S4,11.89)};
\addlegendentry{VAT}
\addplot[fill=colNoPMH, draw=none]
  coordinates {(S1,0.66)(S2,0.42)(S3,0.50)(S4,3.28)};
\addlegendentry{E1 no PMH}
\addplot[fill=colPMH,   draw=none]
  coordinates {(S1,0.66)(S2,0.42)(S3,0.50)(S4,3.34)};
\addlegendentry{E1 (PMH)}
\node[font=\tiny\bfseries, color=colRed, align=right, anchor=north east]
  at (rel axis cs:0.99,0.97) {S4: VAT/PMH\\$3.6\times$};
\end{axis}
\end{tikzpicture}
\caption{Stage-wise drift T07 (ResNet).}
\end{subfigure}
\hfill
\begin{subfigure}[b]{0.31\textwidth}
\begin{tikzpicture}
\begin{axis}[
  pmh bar,
  symbolic x coords={B0,VAT,{E1 no PMH},{E1 (PMH)}},
  xtick=data,
  xticklabel style={font=\footnotesize, rotate=20, anchor=east},
  ylabel={Saliency cosine similarity $\uparrow$},
  ymin=0.44, ymax=0.78,
  width=0.85\linewidth, height=5.5cm,
  bar width=6pt,
  nodes near coords,
  nodes near coords style={font=\tiny, anchor=south},
  point meta=rawy,
  every node near coord/.append style={
    /pgf/number format/fixed,
    /pgf/number format/precision=3},
]
\addplot[fill=colB0, draw=none]
  coordinates {(B0,0.530)(VAT,0.668)({E1 no PMH},0.717)({E1 (PMH)},0.718)};
\end{axis}
\end{tikzpicture}
\caption{Saliency stability T07.}
\end{subfigure}
\caption{\textbf{Cross-domain mechanistic evidence.}
\emph{(a)} Embedding drift at $\sigma{=}0.1$: PMH reduces drift 44\% vs.\ VAT on T01
and 18$\times$ on T02 (graph), the largest reduction of any task.
\emph{(b)} Stage-wise drift on T07: B0/VAT show catastrophic Stage~4 drift ($>$11);
PMH reduces it $3.6\times$.
\emph{(c)} Saliency stability on T07: PMH achieves the highest cosine similarity (0.718),
confirming attentional consistency under perturbation.}
\label{fig:a5}
\end{figure*}


\begin{figure*}[t]
\centering
\begin{subfigure}[b]{0.31\textwidth}
\begin{tikzpicture}
\begin{axis}[
  pmh bar,
  symbolic x coords={T01,T02},
  xtick=data,
  xticklabels={T01 (CIFAR-10), T02 (Graph)},
  xticklabel style={font=\footnotesize, rotate=12, anchor=east},
  ylabel={Embedding drift at $\sigma{=}0.1$},
  ymin=0, ymax=0.82,
  legend pos=north east,
  width=0.88\linewidth, height=5.5cm,
  bar width=4.5pt,
]
\addplot[fill=colB0,   draw=none] coordinates {(T01,0.692)(T02,0.375)};
\addlegendentry{B0}
\addplot[fill=colVAT,  draw=none] coordinates {(T01,0.699)(T02,0.304)};
\addlegendentry{VAT}
\addplot[fill=colPMH,  draw=none] coordinates {(T01,0.385)(T02,0.021)};
\addlegendentry{E1 (PMH)}
\node[font=\tiny, color=colPMH] at (axis cs:T02,0.12) {18$\times$ lower};
\end{axis}
\end{tikzpicture}
\caption{Embedding drift.}
\end{subfigure}
\hfill
\begin{subfigure}[b]{0.34\textwidth}
\begin{tikzpicture}
\begin{axis}[
  pmh bar,
  symbolic x coords={S1,S2,S3,S4},
  xtick=data,
  xlabel={ResNet stage},
  ylabel={Stage drift (lower $=$ better)},
  ymin=0, ymax=14.5,
  legend pos=north west,
  width=0.88\linewidth, height=5.5cm,
  bar width=3.5pt,
]
\addplot[fill=colB0,    draw=none]
  coordinates {(S1,1.82)(S2,1.59)(S3,1.92)(S4,12.68)};
\addlegendentry{B0}
\addplot[fill=colVAT,   draw=none]
  coordinates {(S1,1.67)(S2,0.82)(S3,0.75)(S4,11.89)};
\addlegendentry{VAT}
\addplot[fill=colNoPMH, draw=none]
  coordinates {(S1,0.66)(S2,0.42)(S3,0.50)(S4,3.28)};
\addlegendentry{E1 no PMH}
\addplot[fill=colPMH,   draw=none]
  coordinates {(S1,0.66)(S2,0.42)(S3,0.50)(S4,3.34)};
\addlegendentry{E1 (PMH)}
\node[font=\tiny\bfseries, color=colRed, align=right, anchor=north east]
  at (rel axis cs:0.99,0.97) {S4: VAT/PMH\\$3.6\times$};
\end{axis}
\end{tikzpicture}
\caption{Stage-wise drift T07 (ResNet).}
\end{subfigure}
\hfill
\begin{subfigure}[b]{0.31\textwidth}
\begin{tikzpicture}
\begin{axis}[
  pmh bar,
  symbolic x coords={B0,VAT,{E1 no PMH},{E1 (PMH)}},
  xtick=data,
  xticklabel style={font=\footnotesize, rotate=20, anchor=east},
  ylabel={Saliency cosine similarity $\uparrow$},
  ymin=0.44, ymax=0.78,
  width=0.85\linewidth, height=5.5cm,
  bar width=6pt,
  nodes near coords,
  nodes near coords style={font=\tiny, anchor=south},
  point meta=rawy,
  every node near coord/.append style={
    /pgf/number format/fixed,
    /pgf/number format/precision=3},
]
\addplot[fill=colB0, draw=none]
  coordinates {(B0,0.530)(VAT,0.668)({E1 no PMH},0.717)({E1 (PMH)},0.718)};
\end{axis}
\end{tikzpicture}
\caption{Saliency stability T07.}
\end{subfigure}
\caption{\textbf{Cross-domain mechanistic evidence.}
\emph{(a)} Embedding drift at $\sigma{=}0.1$: PMH reduces drift 44\% vs.\ VAT on T01
and 18$\times$ on T02 (graph), the largest reduction of any task.
\emph{(b)} Stage-wise drift on T07: B0/VAT show catastrophic Stage~4 drift ($>$11);
PMH reduces it $3.6\times$.
\emph{(c)} Saliency stability on T07: PMH achieves the highest cosine similarity (0.718),
confirming attentional consistency under perturbation.}
\label{fig:a5}
\end{figure*}


\begin{figure*}[t]
\centering
\begin{subfigure}[b]{0.31\textwidth}
\begin{tikzpicture}
\begin{axis}[
  pmh bar,
  symbolic x coords={T01,T02},
  xtick=data,
  xticklabels={T01 (CIFAR-10), T02 (Graph)},
  xticklabel style={font=\footnotesize, rotate=12, anchor=east},
  ylabel={Embedding drift at $\sigma{=}0.1$},
  ymin=0, ymax=0.82,
  legend pos=north east,
  width=0.88\linewidth, height=5.5cm,
  bar width=4.5pt,
]
\addplot[fill=colB0,   draw=none] coordinates {(T01,0.692)(T02,0.375)};
\addlegendentry{B0}
\addplot[fill=colVAT,  draw=none] coordinates {(T01,0.699)(T02,0.304)};
\addlegendentry{VAT}
\addplot[fill=colPMH,  draw=none] coordinates {(T01,0.385)(T02,0.021)};
\addlegendentry{E1 (PMH)}
\node[font=\tiny, color=colPMH] at (axis cs:T02,0.12) {18$\times$ lower};
\end{axis}
\end{tikzpicture}
\caption{Embedding drift.}
\end{subfigure}
\hfill
\begin{subfigure}[b]{0.34\textwidth}
\begin{tikzpicture}
\begin{axis}[
  pmh bar,
  symbolic x coords={S1,S2,S3,S4},
  xtick=data,
  xlabel={ResNet stage},
  ylabel={Stage drift (lower $=$ better)},
  ymin=0, ymax=14.5,
  legend pos=north west,
  width=0.88\linewidth, height=5.5cm,
  bar width=3.5pt,
]
\addplot[fill=colB0,    draw=none]
  coordinates {(S1,1.82)(S2,1.59)(S3,1.92)(S4,12.68)};
\addlegendentry{B0}
\addplot[fill=colVAT,   draw=none]
  coordinates {(S1,1.67)(S2,0.82)(S3,0.75)(S4,11.89)};
\addlegendentry{VAT}
\addplot[fill=colNoPMH, draw=none]
  coordinates {(S1,0.66)(S2,0.42)(S3,0.50)(S4,3.28)};
\addlegendentry{E1 no PMH}
\addplot[fill=colPMH,   draw=none]
  coordinates {(S1,0.66)(S2,0.42)(S3,0.50)(S4,3.34)};
\addlegendentry{E1 (PMH)}
\node[font=\tiny\bfseries, color=colRed, align=right, anchor=north east]
  at (rel axis cs:0.99,0.97) {S4: VAT/PMH\\$3.6\times$};
\end{axis}
\end{tikzpicture}
\caption{Stage-wise drift T07 (ResNet).}
\end{subfigure}
\hfill
\begin{subfigure}[b]{0.31\textwidth}
\begin{tikzpicture}
\begin{axis}[
  pmh bar,
  symbolic x coords={B0,VAT,{E1 no PMH},{E1 (PMH)}},
  xtick=data,
  xticklabel style={font=\footnotesize, rotate=20, anchor=east},
  ylabel={Saliency cosine similarity $\uparrow$},
  ymin=0.44, ymax=0.78,
  width=0.85\linewidth, height=5.5cm,
  bar width=6pt,
  nodes near coords,
  nodes near coords style={font=\tiny, anchor=south},
  point meta=rawy,
  every node near coord/.append style={
    /pgf/number format/fixed,
    /pgf/number format/precision=3},
]
\addplot[fill=colB0, draw=none]
  coordinates {(B0,0.530)(VAT,0.668)({E1 no PMH},0.717)({E1 (PMH)},0.718)};
\end{axis}
\end{tikzpicture}
\caption{Saliency stability T07.}
\end{subfigure}
\caption{\textbf{Cross-domain mechanistic evidence.}
\emph{(a)} Embedding drift at $\sigma{=}0.1$: PMH reduces drift 44\% vs.\ VAT on T01
and 18$\times$ on T02 (graph), the largest reduction of any task.
\emph{(b)} Stage-wise drift on T07: B0/VAT show catastrophic Stage~4 drift ($>$11);
PMH reduces it $3.6\times$.
\emph{(c)} Saliency stability on T07: PMH achieves the highest cosine similarity (0.718),
confirming attentional consistency under perturbation.}
\label{fig:a5}
\end{figure*}


\begin{figure*}[t]
\centering
\begin{subfigure}[b]{0.31\textwidth}
\begin{tikzpicture}
\begin{axis}[
  pmh bar,
  symbolic x coords={T01,T02},
  xtick=data,
  xticklabels={T01 (CIFAR-10), T02 (Graph)},
  xticklabel style={font=\footnotesize, rotate=12, anchor=east},
  ylabel={Embedding drift at $\sigma{=}0.1$},
  ymin=0, ymax=0.82,
  legend pos=north east,
  width=0.88\linewidth, height=5.5cm,
  bar width=4.5pt,
]
\addplot[fill=colB0,   draw=none] coordinates {(T01,0.692)(T02,0.375)};
\addlegendentry{B0}
\addplot[fill=colVAT,  draw=none] coordinates {(T01,0.699)(T02,0.304)};
\addlegendentry{VAT}
\addplot[fill=colPMH,  draw=none] coordinates {(T01,0.385)(T02,0.021)};
\addlegendentry{E1 (PMH)}
\node[font=\tiny, color=colPMH] at (axis cs:T02,0.12) {18$\times$ lower};
\end{axis}
\end{tikzpicture}
\caption{Embedding drift.}
\end{subfigure}
\hfill
\begin{subfigure}[b]{0.34\textwidth}
\begin{tikzpicture}
\begin{axis}[
  pmh bar,
  symbolic x coords={S1,S2,S3,S4},
  xtick=data,
  xlabel={ResNet stage},
  ylabel={Stage drift (lower $=$ better)},
  ymin=0, ymax=14.5,
  legend pos=north west,
  width=0.88\linewidth, height=5.5cm,
  bar width=3.5pt,
]
\addplot[fill=colB0,    draw=none]
  coordinates {(S1,1.82)(S2,1.59)(S3,1.92)(S4,12.68)};
\addlegendentry{B0}
\addplot[fill=colVAT,   draw=none]
  coordinates {(S1,1.67)(S2,0.82)(S3,0.75)(S4,11.89)};
\addlegendentry{VAT}
\addplot[fill=colNoPMH, draw=none]
  coordinates {(S1,0.66)(S2,0.42)(S3,0.50)(S4,3.28)};
\addlegendentry{E1 no PMH}
\addplot[fill=colPMH,   draw=none]
  coordinates {(S1,0.66)(S2,0.42)(S3,0.50)(S4,3.34)};
\addlegendentry{E1 (PMH)}
\node[font=\tiny\bfseries, color=colRed, align=right, anchor=north east]
  at (rel axis cs:0.99,0.97) {S4: VAT/PMH\\$3.6\times$};
\end{axis}
\end{tikzpicture}
\caption{Stage-wise drift T07 (ResNet).}
\end{subfigure}
\hfill
\begin{subfigure}[b]{0.31\textwidth}
\begin{tikzpicture}
\begin{axis}[
  pmh bar,
  symbolic x coords={B0,VAT,{E1 no PMH},{E1 (PMH)}},
  xtick=data,
  xticklabel style={font=\footnotesize, rotate=20, anchor=east},
  ylabel={Saliency cosine similarity $\uparrow$},
  ymin=0.44, ymax=0.78,
  width=0.85\linewidth, height=5.5cm,
  bar width=6pt,
  nodes near coords,
  nodes near coords style={font=\tiny, anchor=south},
  point meta=rawy,
  every node near coord/.append style={
    /pgf/number format/fixed,
    /pgf/number format/precision=3},
]
\addplot[fill=colB0, draw=none]
  coordinates {(B0,0.530)(VAT,0.668)({E1 no PMH},0.717)({E1 (PMH)},0.718)};
\end{axis}
\end{tikzpicture}
\caption{Saliency stability T07.}
\end{subfigure}
\caption{\textbf{Cross-domain mechanistic evidence.}
\emph{(a)} Embedding drift at $\sigma{=}0.1$: PMH reduces drift 44\% vs.\ VAT on T01
and 18$\times$ on T02 (graph), the largest reduction of any task.
\emph{(b)} Stage-wise drift on T07: B0/VAT show catastrophic Stage~4 drift ($>$11);
PMH reduces it $3.6\times$.
\emph{(c)} Saliency stability on T07: PMH achieves the highest cosine similarity (0.718),
confirming attentional consistency under perturbation.}
\label{fig:a5}
\end{figure*}
